% Part: computability
% Chapter: recursive-functions
% Section: notations-pr-functions

\documentclass[../../../include/open-logic-section]{subfiles}

\begin{document}

\olfileid{cmp}{rec}{not}
\olsection{بنیادی بازگشت دیاں علامتی لکھتاں}

بنیادی بازگشتی فنکشناں دی ٹھیک استقرائی تعریف
دا اک فائدہ ایہ اے کہ اسی اوہناں نوں
منظم طریقے نال بیان کر سکدے آں۔ مثال لئی،
ہر ایہو جیہے فنکشن نوں اک «علامتی لکھت»
دے سکدے آں، اِس طرح: صفر، جانشین تے
پروجیکشن فنکشناں لئی بالترتیب $\Zero$،
$\Succ$ تے $\Proj{n}{i}$ دیاں علامتاں
ورتو۔ ہن فرض کرو $h$ نوں اک $k$-جگہ
فنکشن~$f$ تے $n$-جگہ فنکشناں $g_0$،
\dots،~$g_{k-1}$ توں ترکیب راہیں
متعین کیتا گیا اے، تے ایہناں فنکشناں نوں
بالترتیب $F$، $G_0$، \dots،~$G_{k-1}$
دیاں علامتی لکھتاں دتیاں گئیاں نیں۔
تاں اک نویں علامت $\fn{Comp}_{k,n}$
ورت کے اسی فنکشن $h$ نوں
$\fn{Comp}_{k,n}[F,G_0,\dots,G_{k-1}]$
نال ظاہر کر سکدے آں۔

بنیادی بازگشت راہیں متعین فنکشناں لئی
وی اِس ورگیاں علامتی لکھتاں ورت سکدے
آں۔ فرض کرو $(k+1)$-جگہ فنکشن~$h$
نوں $k$-جگہ فنکشن~$f$ تے
$(k+2)$-جگہ فنکشن~$g$ توں بنیادی
بازگشت راہیں متعین کیتا گیا اے،
تے $f$ تے~$g$ نوں بالترتیب $F$
تے~$G$ دیاں علامتی لکھتاں دتیاں
گئیاں نیں۔ تاں~$h$ نوں دتی
علامتی لکھت $\fn{Rec}_k[F,G]$ اے۔

یاد کرو کہ جمع دا فنکشن بنیادی بازگشت
راہیں ایویں متعین ہوندا اے:
\begin{align*}
  \Add(x_0, 0) & = \Proj{1}{0}(x_0) = x_0\\
  \Add(x_0, y+1) & = \Succ(\Proj{3}{2}(x_0, y, \Add(x_0, y))) = \Add(x_0, y) +1
\end{align*}
ایتھے~$f$ دا کردار $\Proj{1}{0}$
ادا کردا اے، تے~$g$ دا کردار
$\Succ(\Proj{3}{2}(x_0, y, z))$
ادا کردا اے۔ اِس نوں
$\fn{Comp}_{1,3}[\Succ,\Proj{3}{2}]$
دی علامتی لکھت دتی جاندی اے، کیونکہ
ایہ $1$-جگہ فنکشن~$\Succ$ تے
$3$-جگہ فنکشن~$\Proj{3}{2}$ توں
ترکیب راہیں فنکشن متعین کرن دا نتیجہ
اے۔ اِس ترتیب نال اسی جمع دے فنکشن
نوں ایویں ظاہر کر سکدے آں:
\[
\fn{Rec}_1[\Proj{1}{0},\fn{Comp}_{1,3}[\Succ,\Proj{3}{2}]].
\]
ایہ علامتی لکھتاں بعض ویلے کم آؤندیاں
نیں، مثلاً بنیادی بازگشتی فنکشناں نوں
گِن کے ترتیب وار درج کرن ویلے۔

\begin{prob}
  $\Mult$ لئی پوری بنیادی بازگشتی علامتی لکھت دیو۔
\end{prob}

\end{document}
